% ============================================================
% Local Spillovers from Recreational Marijuana Legalization:
% Housing Markets and Drug-Overdose Mortality Near State Borders
% Formatted to Economic Inquiry (Wiley/WEAI) submission requirements.
% Compile: pdfLaTeX -> BibTeX -> pdfLaTeX -> pdfLaTeX (Overleaf: Recompile)
% ============================================================
\documentclass[11pt]{article}

\usepackage[margin=1in]{geometry}
\usepackage{setspace}
\usepackage{amsmath,amssymb}
\usepackage{graphicx}
\usepackage{booktabs}
\usepackage{caption}
\usepackage[authoryear,round]{natbib}
\usepackage[hidelinks]{hyperref}
\usepackage{enumitem}
\usepackage{titlesec}

\doublespacing
\titleformat{\section}{\normalfont\bfseries\large}{\thesection.}{0.5em}{}
\titleformat{\subsection}{\normalfont\bfseries}{\thesubsection}{0.5em}{}

\title{}
\author{}
\date{}

\begin{document}

% ============================================================
% TITLE PAGE
% ============================================================
\begin{titlepage}
\begin{center}
\vspace*{1.5cm}
{\LARGE \bfseries Local Spillovers from Recreational Marijuana\\
Legalization: Housing Markets and Drug-Overdose\\
Mortality Near State Borders\par}
\vspace{0.3cm}
{\large \itshape Evidence from a Continuous-Distance Border Design, 2020--2023\par}
\vspace{1.2cm}
{\large Author name(s) withheld for review\par}
\vspace{0.3cm}
{\normalsize Working Paper --- July 2026\par}
\end{center}

\vspace{1.2cm}
\noindent\textbf{Abstract} \\[4pt]
\noindent We test for local economic and public-health spillovers from
proximity to a newly legal recreational-marijuana market, adapting a
switcher/control county design from a companion project to a tighter
50-mile radius: counties in never-legalizing states whose distance to
the nearest legal border fell below 50 miles between 2020 and 2023
(``switchers'') versus counties that remained farther away
(``controls''). We examine two outcomes: county home values and rents
(Zillow), and drug-overdose mortality (CDC). Home-value growth is
modestly higher in counties that stayed farther from a legal border
($p=0.073$), a possible disamenity signal. Switcher counties show a
significantly larger overdose-mortality decline than controls
($p=0.030$), consistent with an opioid-substitution channel, though
this does not survive a continuous-distance specification. We
interpret both results as suggestive rather than conclusive. \\[10pt]
\noindent\textbf{JEL Classification:} R21, I18, I12, K42 \\
\noindent\textbf{Keywords:} recreational marijuana; housing prices;
drug overdose mortality; opioid substitution; border spillovers;
difference-in-differences

\vfill
\noindent\textbf{Abbreviations} \\[2pt]
\noindent
\begin{tabular}{@{}ll@{}}
DiD    & Difference-in-differences \\
FIPS   & Federal Information Processing Standard (county code) \\
NCHS   & National Center for Health Statistics \\
SE     & Standard error \\
VSRR   & Vital Statistics Rapid Release \\
ZHVI   & Zillow Home Value Index \\
ZORI   & Zillow Observed Rent Index \\
\end{tabular}

\vspace{0.5cm}
\noindent{\footnotesize A separate Disclosure Statement accompanies this
submission per Economic Inquiry policy and is not reproduced in this file.}
\end{titlepage}

% ============================================================
% 1. INTRODUCTION
% ============================================================
\section{Introduction}

Does proximity to a legal recreational-marijuana market change local
economic and public-health outcomes for residents who do not
themselves cross the border? A companion project examines this
question through the lens of drug-enforcement substitution, using a
switcher/control county design -- counties in never-legalizing states
whose distance to the nearest legal-market border fell below 100 miles
between 2020 and 2023 -- and finds no robust evidence that marijuana
and alcohol arrests move as substitutes. Arrest data, however, are
enforcement-mediated rather than purely behavioral, and enforcement
intensity may respond to legalization independently of underlying
behavior. This paper adapts the same treatment design, at a tighter
50-mile switcher/control radius (Section~2.1), to examine two outcomes
that are not enforcement-mediated in the same way: local housing
markets and drug-overdose mortality.

The housing question follows a capitalization logic: if proximity to a
legal market is an amenity (convenience, local retail and tourism
spillovers) or a disamenity (traffic, stigma, perceived crime risk),
forward-looking home buyers and renters should price it in.
\citet{ChengMayerMayer2018} find that legalizing retail marijuana
sales in Colorado raised nearby home values, but their design compares
locations \emph{within} a legal state to dispensary proximity, not
proximity to an external legal-market \emph{border} from the
non-legal side -- the mechanism and the sign are not obviously the
same across these two settings.

The overdose-mortality question follows a substitution logic
distinct from the alcohol-arrest test in the companion paper: if
marijuana substitutes for opioids in pain management or recreational
use, easier access should reduce opioid-involved mortality nearby.
\citet{Bachhuber2014} found early evidence of this using state medical-
marijuana laws, but \citet{Shover2019} show the association reverses
sign once more years of data are included, and \citet{PowellPaculaJacobson2018}
find the effect is specific to certain legal designs (dispensary
access) rather than medical-marijuana laws generally. Whether a
\emph{recreational}-legalization, cross-border-proximity design shows
any such pattern in the 2020--2023 wave is, to our knowledge,
untested.

We find a marginally significant increase in home-value growth for
counties that remained farther from a legal border ($p=0.073$ in the
primary continuous-distance specification), not significant in the
binary switcher specification. We find a significant decline in
drug-overdose mortality for switcher counties relative to controls
($p=0.030$), but this does not survive the continuous-distance
specification, and rests on a sample reduced by roughly 60\% due to
small-county data suppression. We report both results with their
specific caveats rather than treating either as a settled finding.

% ============================================================
% 2. DATA
% ============================================================
\section{Data}

\subsection{Treatment Design (Adapted from the Companion Project, at a Tighter Radius)}

We reuse the underlying county panel and distance computation from the
companion marijuana-alcohol substitution project -- 32 states that had
not legalized recreational marijuana through 2023, with each county's
distance to the nearest legal-market border computed in miles for 2020
and 2023 (Haversine distance between county centroids) -- but redefine
the switcher/control cutoff at a tighter 50-mile radius rather than the
companion project's 100 miles. We first attempted a 20-mile radius, but
this leaves only 32 switcher counties nationally, too few once split
across three outcomes with very different data coverage (Section~2.2--2.3);
50 miles yields 212 switcher counties nationally and is the radius used
throughout this paper. A county is a ``switcher'' if its distance to
the nearest legal-market border was at least 50 miles in 2020 and fell
below 50 miles by 2023; a ``control'' if it remained at least 50 miles
away in both years; counties within 50 miles of a legal border in
\emph{both} years (54 counties) are excluded from the main sample, as
in the companion paper.

The primary treatment variable is the raw change in distance,
\begin{equation}
  \Delta\text{Dist}_i \;=\; \text{Dist}_{i,2023} \;-\; \text{Dist}_{i,2020},
  \label{eq:distmi}
\end{equation}
negative for counties that became closer to a newly-legal border. We
estimate first-difference regressions
\begin{equation}
  \Delta Y_i \;=\; \alpha \;+\; \beta \cdot \frac{\Delta\text{Dist}_i}{50} \;+\; \varepsilon_i,
  \label{eq:distreg}
\end{equation}
scaling into 50-mile units to match the switcher radius, alongside the
binary switcher indicator and a $\Delta\log(\text{distance})$
robustness check, with standard errors clustered by state throughout
\citep{CameronMiller2015, BertrandDufloMullainathan2004}. The full
panel available to this treatment design is 2,347 counties (212
switchers, 2,081 controls, 54 always-near); the two outcome sources
below cover different, smaller subsets of it.

\subsection{Housing: Zillow ZHVI and ZORI}

We use Zillow Research's publicly available, county-level Zillow Home
Value Index (ZHVI, a smoothed, seasonally-adjusted measure of the
typical home value for the middle price tier) and Zillow Observed
Rent Index (ZORI), both measured as of December 2020 and December
2023. ZHVI covers 3,071 counties nationally, 2,191 of them in our
main treatment sample (207 switchers, 1,984 controls). ZORI has
substantially sparser rural coverage -- 1,360 counties nationally, only
361 in our main sample (43 switchers, 318 controls) -- because Zillow
requires a minimum volume of active rental listings to compute a
county-level rent index, which many small, non-metropolitan counties
in our sample do not meet. We use the log change in each index,
$\Delta \log(\text{ZHVI})$ and $\Delta \log(\text{ZORI})$, as the
outcome.

\subsection{Drug-Overdose Mortality: CDC VSRR Provisional County Data}

County-level, cause-specific mortality is not available through any
open, unauthenticated API: CDC WONDER, the relevant query tool, is not
reachable from our data-collection environment, and its underlying
Multiple/Underlying Cause of Death files require a manual query
through its web interface. We instead use the CDC's Vital Statistics
Rapid Release (VSRR) ``Provisional County-Level Drug Overdose Death
Counts,'' published on the open data.cdc.gov Socrata platform, which
reports 12-month-ending drug-overdose death counts by county of
residence. We take the count for the 12 months ending December 2020
and December 2023 as our two periods, divide by county population to
get a rate per 100,000, and take the first difference as the outcome.

Two limitations of this source are consequential enough to state
plainly. First, per NCHS confidentiality standards, any county-month
cell with 1--9 deaths is suppressed; this removes 1,505 of 3,144
counties nationally in 2020 and 1,425 in 2023, and correspondingly
shrinks our main-sample overdose regression to 928 counties (81
switchers, 847 controls) -- less than half of the full treatment
panel, disproportionately excluding the smallest, most rural counties,
and especially costly at a 50-mile radius since the switcher pool is
already small. Second, this is \emph{provisional} surveillance data
subject to revision as death-certificate investigations are completed
and jurisdictions finalize reporting (the data file's own completeness
field averages well under 100\% for our sample period); it is not the
final vital-statistics record. We treat the overdose result below as
suggestive for these reasons.

% ============================================================
% 3. RESULTS
% ============================================================
\section{Results}

\begin{table}[htbp]
\centering
\caption{Effect on Home Values, Rent, and Overdose Mortality}
\label{tab:regD}
\begin{tabular}{@{}lccc@{}}
\toprule
 & Switcher (binary) & Distance, per 50 mi & Log-distance (robustness) \\
\midrule
Home value (log) & $+0.011$ (0.018) & $+0.009^{*}$ (0.005) & $+0.015$ (0.010) \\
 & $p=0.549$ & $p=0.073$ & $p=0.139$ \\
\addlinespace
Rent (log) & $-0.014$ (0.012) & $-0.002$ (0.003) & $+0.010$ (0.009) \\
 & $p=0.238$ & $p=0.586$ & $p=0.257$ \\
\addlinespace
Overdose rate per 100k & $-3.546^{**}$ (1.632) & $-0.554$ (0.340) & $+0.738$ (1.292) \\
 & $p=0.030$ & $p=0.103$ & $p=0.568$ \\
\bottomrule
\end{tabular}
\vspace{2pt}
\begin{minipage}{0.95\textwidth}
\footnotesize Note: Outcome is the change (2020--2023) in the log
index (home value, rent) or arrest-rate-style count per 100{,}000
(overdose). SE clustered by state in parentheses. $N=2{,}191$ (home
value), $N=361$ (rent), $N=928$ (overdose) -- see Section~2 for why
coverage differs by outcome. $^{*}p<0.10$, $^{**}p<0.05$.
\end{minipage}
\end{table}

\begin{figure}[htbp]
  \centering
  \includegraphics[width=0.95\textwidth]{figures/figD1_coefplot.png}
  \caption{Effect on Home Values, Rent, and Overdose Mortality, All Three Specifications (95\% CI).}
  \label{fig:coefD}
\end{figure}

\subsection{Housing Markets}

The primary, continuous-distance specification finds that a 50-mile
\emph{increase} in distance to the nearest legal border is associated
with 0.9 log points \emph{more} home-value growth ($p=0.073$) --
equivalently, counties that became 50 miles \emph{closer} to a legal
border saw about 0.9 log points \emph{less} home-value growth over
2020--2023 than counties that stayed where they were. This is only
marginally significant and does not appear in the binary switcher
specification ($p=0.549$) or the log-distance robustness check
($p=0.139$), so we read it as weak evidence, at best, of a local
disamenity effect on home prices -- the opposite sign from
\citet{ChengMayerMayer2018}'s within-state, dispensary-proximity
result, though the settings are different enough (cross-border
proximity from the non-legal side, versus within-state proximity to
retail outlets) that the two findings are not in direct tension.
Rent shows no significant relationship with any specification, though
the rent sample is small enough (361 counties, only 43 of them
switchers at this tighter radius) that this is likely underpowered
rather than a precise null (Figure~\ref{fig:coefD}).

\begin{figure}[htbp]
  \centering
  \includegraphics[width=0.95\textwidth]{figures/figD2_treatment_scatter.png}
  \caption{The Continuous-Treatment Regression: Distance Change vs. Outcome Change, Home Value and Overdose Rate.}
  \label{fig:scatterD}
\end{figure}

\subsection{Drug-Overdose Mortality}

Switcher counties saw drug-overdose death rates fall by 3.5 more
per 100,000 than control counties over the same period ($p=0.030$),
a result of a size worth taking seriously (Figure~\ref{fig:overdosetrend}
shows both groups' raw trends). However, the continuous-distance
specification -- which uses the same underlying variation and is our
preferred specification in the companion paper -- is not significant
($p=0.103$, and the log-distance robustness check flips sign entirely,
$p=0.568$). A result present in the binary comparison but not the
continuous one is a meaningfully weaker form of evidence than a result
that holds across specifications, and combined with the sample
reduction from suppression (Section~2.3), we do not treat this as a
confirmed opioid-substitution effect -- only as a pattern worth further
investigation, ideally once final (non-provisional) vital-statistics
data are available for these years.

\begin{figure}[htbp]
  \centering
  \includegraphics[width=0.7\textwidth]{figures/figD3_overdose_trends.png}
  \caption{Drug Overdose Death Rate, Switchers vs. Controls, 2020--2023.}
  \label{fig:overdosetrend}
\end{figure}

% ============================================================
% 4. DISCUSSION
% ============================================================
\section{Discussion}

Taken together with the companion paper's null substitution result on
arrest data, this paper's two results sit at different points on a
spectrum of confidence. The housing result is directionally consistent
across two of three specifications but only marginally significant in
the one we most trust, and small in magnitude. The overdose result is
larger and more significant in its binary form, but fails to survive
the continuous-distance specification that is otherwise our preferred
test throughout this line of work -- exactly the kind of
specification-dependence that \citet{Shover2019} caution is common in
this literature once enough data years are examined. Neither result
should be read as a settled finding; both are more useful as a map of
where a dedicated follow-up (a longer panel, final rather than
provisional mortality data, or a larger rent sample) would have the
best chance of a conclusive answer.

% ============================================================
% 5. LIMITATIONS
% ============================================================
\section{Limitations}

\begin{enumerate}[label=(\arabic*), leftmargin=*]
\item \textbf{Overdose mortality data are provisional, not final.}
CDC's VSRR county-level file is explicitly a rapid-release product
subject to revision as death investigations close; a final-data
replication is not yet possible for 2023.
\item \textbf{Small-cell suppression removes the majority of small
counties from the overdose sample.} Only 928 of 2,347 treatment-panel
counties have non-suppressed data in both years, and the excluded
counties are disproportionately rural -- exactly the counties where a
border-proximity design has the most identifying variation, since they
sit closest to state lines.
\item \textbf{Rent (ZORI) coverage is sparse for the same reason.}
Zillow requires a minimum active-listing volume to compute a county
rent index; only 361 of 2,347 counties qualify, limiting that
regression's power.
\item \textbf{The 50-mile switcher radius is itself a data-driven
compromise.} We could not use the originally requested 20-mile radius
because it leaves only 32 switcher counties nationally; 50 miles was
the next threshold tested and is not derived from any theoretical
ideal catchment distance for either outcome.
\item \textbf{Two-period design.} As in the companion paper, the
switcher/control comparison uses only 2020 and 2023, precluding a
formal pre-trend test.
\item \textbf{No demand-side decomposition.} We cannot distinguish,
within the housing result, whether any disamenity effect reflects
traffic, crime perception, stigma, or some other specific channel; nor,
within the overdose result, whether any effect reflects substitution
in pain management, recreational substitution, or a compositional
change in who lives in switcher counties.
\end{enumerate}

% ============================================================
% 6. CONCLUSION
% ============================================================
\section{Conclusion}

Adapting the switcher/control, continuous-distance treatment design
from our companion marijuana-alcohol substitution paper to a tighter
50-mile radius, we find
suggestive but not conclusive evidence of two local spillovers from
proximity to a newly legal recreational-marijuana market: modestly
slower home-value growth in counties that grew closer to a legal
border, and a larger decline in drug-overdose mortality in switcher
counties specifically. Both results are sensitive to specification or
data-quality caveats detailed above, and we recommend they be treated
as motivation for a more targeted follow-up -- particularly once final
(non-provisional) mortality data for 2023 are released -- rather than
as standalone evidence of either a housing disamenity or an opioid-
substitution effect.

% ============================================================
% REFERENCES
% ============================================================
\bibliographystyle{apalike}
\bibliography{references}

\end{document}
